\documentclass{rapportECL}
\usepackage{listings}
\usepackage{float}
\usepackage{xcolor}
\usepackage{csquotes}
\usepackage{booktabs}
\usepackage{minted}
\usepackage{caption}
\usepackage{amssymb}
\usepackage[linesnumbered,ruled,vlined]{algorithm2e}
\usepackage{placeins}
\usepackage[french]{babel}
\usepackage{url}
\usepackage{graphicx}
\usepackage{array}
\usepackage{colortbl}
\usepackage{tikz}
\usetikzlibrary{shapes.geometric, arrows.meta, positioning, bending}
\usepackage{cite}
\usepackage{hyperref}
\usepackage{pgfplots}
\pgfplotsset{compat=1.18}

\hypersetup{
    colorlinks=true,
    linkcolor=blue,
    citecolor=green,
    filecolor=magenta,
    urlcolor=cyan,
}

% --- Infos rapport ---
\title{RAG-Tutor}
\definecolor{codebg}{RGB}{248,250,252}
\definecolor{codeframe}{RGB}{203,213,225}

\lstdefinestyle{ragcode}{
    language=Python,
    backgroundcolor=\color{codebg},
    basicstyle=\ttfamily\small,
    keywordstyle=\color{blue}\bfseries,
    commentstyle=\color{green!60!black},
    stringstyle=\color{orange!80!black},
    frame=single,
    framerule=0.6pt,
    rulecolor=\color{codeframe},
    numbers=left,
    numberstyle=\footnotesize\color{gray},
    stepnumber=1,
    numbersep=10pt,
    breaklines=true,
    showspaces=false,
    showstringspaces=false,
    tabsize=4,
    captionpos=b,
}
\lstset{style=ragcode}

\begin{document}

% ----------- Page de garde ---------
\titre{RAG-Tutor : Assistant Pédagogique Intelligent\\par Retrieval-Augmented Generation}
\UE{\textbf{TER — Travail d'Étude et de Recherche}}
\enseignant{À COMPLÉTER — Nom du tuteur}
\eleves{À COMPLÉTER — Ton nom}

\fairemarges
\fairepagedegarde
\clearpage

% ----------- Résumé (FR) ---------
\section*{Résumé}
\addcontentsline{toc}{section}{Résumé}

RAG-Tutor est un assistant pédagogique destiné aux étudiants en Deep Learning, capable de répondre à des questions en s'appuyant exclusivement sur le corpus du cours \emph{Dive into Deep Learning} (d2l.ai). Le système repose sur une architecture RAG (Retrieval-Augmented Generation) associant un retriever hybride (BM25 + dense BGE-M3 + cross-encoder) à un modèle de langage local (qwen2.5:14b via Ollama). Un mécanisme de refus à deux niveaux — score du cross-encoder en pré-génération et juge LLM en post-génération — garantit que le système ne répond que lorsque l'information est présente dans le corpus. Évalué sur un dataset de 45 questions (30 answerable, 15 unanswerable hors-domaine), le système atteint un hit@4 de 0,967, un MRR de 0,867, et une exactitude de refus de 100\% sur les questions hors-corpus. Une étude d'ablation en quatre étapes valide la contribution de chaque composant.

\bigskip
\noindent\textbf{Mots-clés :}
Retrieval-Augmented Generation (RAG),
Recherche sémantique,
Retriever hybride,
Cross-encoder,
Mécanisme de refus,
Modèles de langage (LLM),
Ollama,
Évaluation automatique,
Ablation.

% ----------- Abstract (EN) ---------
\section*{Abstract}
\addcontentsline{toc}{section}{Abstract}

RAG-Tutor is an educational assistant for Deep Learning students, designed to answer questions based solely on the \emph{Dive into Deep Learning} (d2l.ai) course corpus. The system uses a Retrieval-Augmented Generation (RAG) architecture combining a hybrid retriever (BM25 + BGE-M3 dense + cross-encoder) with a local language model (qwen2.5:14b via Ollama). A two-level refusal mechanism — cross-encoder score pre-generation and LLM judge post-generation — ensures the system only answers when information is present in the corpus. Evaluated on a 45-question dataset (30 answerable, 15 out-of-domain unanswerable), the system achieves hit@4 of 0.967, MRR of 0.867, and 100\% refusal accuracy on out-of-domain questions. A four-step ablation study validates each component's contribution.

\bigskip
\noindent\textbf{Keywords:}
Retrieval-Augmented Generation (RAG),
Semantic Search,
Hybrid Retriever,
Cross-encoder,
Refusal Mechanism,
Large Language Models (LLM),
Ollama,
Automatic Evaluation,
Ablation Study.

\clearpage

% ----------- Remerciements ---------
\newpage
\section*{Remerciements}
\addcontentsline{toc}{section}{Remerciements}

À COMPLÉTER

\newpage

% ----------- Table des matières ---------
\addcontentsline{toc}{section}{Table des matières}
\tableofcontents
\newpage
\addcontentsline{toc}{section}{Table des figures}
\listoffigures
\newpage

% ================================================================
% CHAPITRE 1 — INTRODUCTION
% ================================================================
\section{Introduction}

\subsection{Contexte et problématique}

À COMPLÉTER — Contexte : assistants pédagogiques IA, problématique des hallucinations, objectif : un tuteur fiable basé sur d2l.ai.

\subsection{Objectifs du projet}

À COMPLÉTER

\subsection{Organisation du rapport}

À COMPLÉTER

% ================================================================
% CHAPITRE 2 — ÉTAT DE L'ART
% ================================================================
\section{État de l'art et fondements théoriques}

\subsection{Retrieval-Augmented Generation (RAG)}

À COMPLÉTER

\subsection{Retrieval : approches denses, parcimonieuses et hybrides}

\subsubsection{BM25 (retrieval parcimonieux)}
\subsubsection{Embeddings denses (BGE-M3)}
\subsubsection{Fusion hybride : Reciprocal Rank Fusion (RRF)}
\subsubsection{Reranking : cross-encoder}

\subsection{Mécanismes de refus (refusal gate)}

\subsection{Modèles de langage locaux (Ollama)}

\subsection{Synthèse et positionnement}

% ================================================================
% CHAPITRE 3 — CONCEPTION DU SYSTÈME
% ================================================================
\section{Conception du système}

\subsection{Architecture globale}

% Schéma du pipeline ici
\begin{figure}[H]
\centering
% \includegraphics[width=\textwidth]{images/architecture_pipeline.png}
\caption{Architecture du pipeline RAG-Tutor}
\label{fig:architecture}
\end{figure}

\subsection{Chunking et indexation}

\subsubsection{Stratégie parent-child}
\subsubsection{Corpus et prétraitement}
\subsubsection{Base vectorielle ChromaDB}

\subsection{Query processing}

\subsection{Pipeline de retrieval}

\subsubsection{Mode dense (baseline)}
\subsubsection{Mode hybride (BM25 + dense)}
\subsubsection{Mode hybride avec reranking (cross-encoder)}

\subsection{Génération de la réponse}

\subsubsection{System prompt}
\subsubsection{Modèle de génération (qwen2.5:14b)}

\subsection{Mécanismes de refus}

\subsubsection{Mécanisme 1 : score du reranker (pré-génération)}
\subsubsection{Mécanisme 2 : juge LLM (post-génération)}

% ================================================================
% CHAPITRE 4 — DATASET D'ÉVALUATION
% ================================================================
\section{Dataset d'évaluation}

\subsection{Génération du dataset}

\subsection{Catégories de questions}

\subsubsection{Single passage}
\subsubsection{Multi passage}
\subsubsection{Unanswerable}

\subsection{Correction du biais unanswerable}

% ================================================================
% CHAPITRE 5 — IMPLÉMENTATION
% ================================================================
\section{Implémentation}

\subsection{Stack technique}

\subsection{Ingestion et chunking}

\subsection{Retriever hybride}

% Extrait de code : retriever_hybride.py
\begin{lstlisting}[caption={Fusion RRF pondérée}, label={lst:rrf}]
# retriever_hybride.py — RRF pondéré
RRF_ALPHA = 0.6  # avantage au dense
def _rrf_fusion(dense_hits, bm25_hits, k=60, alpha=RRF_ALPHA):
    scores = defaultdict(float)
    for rank, hit in enumerate(dense_hits, start=1):
        scores[hit["meta"]["parent_id"]] += alpha / (k + rank)
    for rank, hit in enumerate(bm25_hits, start=1):
        scores[hit["meta"]["parent_id"]] += (1 - alpha) / (k + rank)
    ...
\end{lstlisting}

\subsection{Reranker cross-encoder}

\subsection{Mécanismes de refus}

% Extrait de code : should_refuse_reranker
\begin{lstlisting}[caption={Refus basé sur le score du cross-encoder}, label={lst:refusal}]
RERANKER_REFUSAL_THRESHOLD = 0.1119

def should_refuse_reranker(hits, threshold=RERANKER_REFUSAL_THRESHOLD):
    if not hits:
        return True
    top_score = hits[0].get("dist")
    if top_score is None:
        return False
    return top_score < threshold
\end{lstlisting}

\subsection{Pipeline complet}

% Extrait de code : pipeline.py (answer)
\begin{lstlisting}[caption={Fonction principale du pipeline}, label={lst:pipeline}]
def answer(query, k=4, use_query_processing=True, use_refusal_gate=True):
    # 1. Query processing
    if use_query_processing:
        rewritten, sub_queries = process_query(query)
    ...
    # 2. Retrieval (round-robin si multi-query)
    all_hits = retrieve_and_merge(queries, k)
    # 3. M1: refus pré-génération (reranker)
    if use_refusal_gate and should_refuse_reranker(all_hits):
        return refuse()
    # 4. Génération
    answer_text = generate(query, all_hits)
    # 5. M2: refus post-génération (LLM judge)
    if use_refusal_gate and not verify_answer(query, answer_text, all_hits):
        return refuse()
    return answer_text
\end{lstlisting}

% ================================================================
% CHAPITRE 6 — PROTOCOLE EXPÉRIMENTAL ET RÉSULTATS
% ================================================================
\section{Protocole expérimental et résultats}

\subsection{Métriques d'évaluation}

\begin{itemize}
    \item \textbf{hit@k} : le document pertinent est-il dans le top-k ?
    \item \textbf{MRR} : à quel rang apparaît le premier document pertinent ?
    \item \textbf{nDCG@k} : pertinence graduée du top-k
    \item \textbf{Precision@k} : proportion de documents pertinents dans le top-k
    \item \textbf{Refusal correctness} : le système refuse-t-il correctement ?
    \item \textbf{Refusal Precision / Recall / F1}
\end{itemize}

\subsection{Intervalles de confiance bootstrap}

\subsection{Étude d'ablation}

\begin{table}[H]
\centering
\caption{Résultats de l'étude d'ablation (test\_set.json, 45 questions)}
\label{tab:ablation}
\begin{tabular}{lcccc}
\toprule
\textbf{Configuration} & \textbf{hit@4} & \textbf{MRR} & \textbf{nDCG@4} & \textbf{Refus} \\
\midrule
Dense (baseline)      & 0,900 & 0,811 & 0,942 & — \\
Dense + QP corrigé     & 0,900 & 0,800 & 0,932 & — \\
Hybrid (BM25 + dense)  & 0,900 & 0,772 & 0,949 & — \\
\textbf{Hybrid + Reranker} & \textbf{0,967} & \textbf{0,867} & \textbf{0,956} & \textbf{1,000} \\
\bottomrule
\end{tabular}
\end{table}

\subsection{Résultats par catégorie}

\subsection{Calibration du seuil de refus}

\begin{figure}[H]
\centering
% \includegraphics[width=0.8\textwidth]{images/calibration_reranker.png}
\caption{Distribution des scores du cross-encoder : answerable vs unanswerable (zéro overlap)}
\label{fig:calibration}
\end{figure}

\subsection{Analyse du bug Query Processing}

% ================================================================
% CHAPITRE 7 — DISCUSSION
% ================================================================
\section{Discussion}

\subsection{Forces du système}

\subsection{Limites}

\subsection{Comparaison avec l'état de l'art}

% ================================================================
% CHAPITRE 8 — CONCLUSION ET PERSPECTIVES
% ================================================================
\section{Conclusion et perspectives}

\subsection{Bilan du projet}

\subsection{Améliorations techniques possibles}

\subsection{Ouverture : déploiement et usage en conditions réelles}

% ================================================================
% BIBLIOGRAPHIE
% ================================================================
\newpage
\bibliographystyle{plain}
\bibliography{references}

\end{document}
